\chapter{总结与展望}
\label{cha:conclusion}

\section{全文总结}

本文围绕企业私有文档知识服务场景，完成了一套基于检索增强生成的智能知识管理系统的需求分析、系统设计、关键实现与系统测试。论文试图解决的，并不是单个问答模型是否足够强，而是企业内部文档如何在权限约束下进入统一知识链路，并进一步支撑检索、问答和文件消费等连续业务过程。从这个意义上说，本文对应的不是单一算法实验，而是一个面向真实使用边界的工程系统。

在需求分析与系统设计阶段，本文把组织边界、文档归属、公开属性、状态恢复和来源可追溯性一并纳入系统建设范围。围绕这些约束，系统建立了由用户角色、组织标签、私人空间和公开资源共同构成的权限框架，并据此组织文档接入、知识检索和结果消费流程。这使系统实现不再停留于“模型接上知识库”的局部演示，而转向对访问边界和文档生命周期的整体处理。

在工程实现上，系统形成了较为完整的业务闭环。后端以前后端分离架构为基础，实现了基于 JWT 的身份认证、基于组织标签的资源隔离、面向大文件场景的分片上传与断点续传、基于 Kafka 的异步解析与入库流程，以及结合 Elasticsearch 的混合检索能力；在交互侧，系统通过 WebSocket 提供流式回答，并维护会话历史、来源编号与文件标识之间的关联关系。本文的实现结果说明，企业私有文档场景下的 RAG 系统，其可用性首先取决于多模块能否稳定协同，而不是取决于某一个局部环节是否表现突出。

第五章的测试进一步给出了这一结论的实际边界。围绕第二章提出的功能需求与非功能需求，系统对身份建立、文档接入、检索权限、文件服务和流式交互进行了系统级验证。结果表明，500 份测试文档均完成上传、续传与对象合并，100 份全链路样本中有 99 份完成解析与索引写入；40 条检索查询全部在 Top5 内命中目标文档，25 组权限矩阵检查未出现越权返回，37 项文件服务检查全部通过。与此同时，流式回答与历史记录读取能力已经成立，但多轮历史追问仍存在语义偏移。由此可以看出，当前版本已经形成主要业务闭环，但在入库稳定性和会话连续性方面仍保留着明确的改进空间。

总体而言，本文的主要工作不在于提出一种新的生成算法，而在于把 RAG 能力约束进一个具有权限控制、状态管理和文档生命周期支撑的软件系统中，并给出可核查的工程验证结果。这一工作说明，企业私有文档场景下的知识问答系统可以沿着文档接入、权限治理、混合检索和流式交互的路径落实为可运行的软件形态。

\section{后续展望}

尽管本文已经将企业私有文档场景中的核心工程链路基本打通，但第五章的测试结果也说明，系统距离长期稳定可用还有一段路要走。

当前最需要继续打磨的，仍然是检索结果的排序质量。现有混合检索已经能够在多数查询上召回目标文档，这说明召回链路本身基本成立；真正直接影响问答体验的，是相关文档能否稳定排到前列，尤其是首位结果是否足够可信。后续工作应更多围绕查询改写、重排序策略、证据片段组织以及人工标注评测集的补充展开，逐步减少生成阶段对提示词修补和经验性调参的依赖。

如果系统面向更大的文档规模和用户规模运行，问题又会转到运维侧。异步入库、断点续传和权限过滤只是基础能力，真正进入业务环境后，还需要有更细的任务监控、失败重试、索引增量更新、审计日志与性能观测机制。很多时候，企业系统的瓶颈并不出现在一次问答是否生成成功，而是出现在链路长时间运行后是否仍然可追踪、可恢复、可维护。这部分工作不如模型效果那样直观，却往往决定系统能否真正进入日常业务流程。

另外，当前知识处理对象仍以可解析文本为主，这一前提在本轮实验范围内是成立的，但放到真实企业资料中并不充分。制度文件、技术规范、扫描件、图表混排文档，甚至版式复杂的项目材料，都会对解析、切分和引用追溯提出更高要求。后续若要继续扩大系统适用范围，就需要在保持权限约束和来源可追溯性的前提下，引入更丰富的文档理解与结构化表示能力。

总的来看，本文完成的更像是一个边界清楚的阶段性成果。它证明了企业私有文档场景下的 RAG 系统可以沿着文档接入、权限控制、混合检索和流式问答这条链路落实到实际工程中；至于更稳定的排序效果、更复杂的文档理解以及更成熟的运维支撑，还需要在后续工作中继续推进。
